\documentclass[10pt,a4paper]{amsart}
\usepackage{amsmath,amssymb,amsthm,graphicx}
% Retained minimal preamble for the Stacks Project chapter more-algebra.tex (More on Algebra).
% The chapter file pinned-original-source.tex carries only \input{preamble} on its first line, so no
% preamble text lives inside the chapter itself.  The Stacks Project uses one single shared preamble for
% every chapter; the snapshot stages it identically as sources/stacks-chow/preamble.tex and
% sources/stacks-etale-cohomology/preamble.tex.  The block below is copied mechanically, line for line,
% from sources/stacks-chow/preamble.tex lines 14-18 (the Xy-pic setup) and 154-217 (the native theorem
% declarations together with the mathematical macros).  Nothing was added, removed or redefined; the
% wrapper class and the amsmath/amssymb/amsthm/graphicx base packages are supplied by the assembler.
% Chapter: more-algebra.tex -- More on Algebra.
% For commutative diagrams we use Xy-pic
\usepackage[all]{xy}

% We use 2cell for 2-commutative diagrams.
\xyoption{2cell}

% Theorem environments.
%
\theoremstyle{plain}
\newtheorem{theorem}[subsection]{Theorem}
\newtheorem{proposition}[subsection]{Proposition}
\newtheorem{lemma}[subsection]{Lemma}

\theoremstyle{definition}
\newtheorem{definition}[subsection]{Definition}
\newtheorem{example}[subsection]{Example}
\newtheorem{exercise}[subsection]{Exercise}
\newtheorem{situation}[subsection]{Situation}

\theoremstyle{remark}
\newtheorem{remark}[subsection]{Remark}
\newtheorem{remarks}[subsection]{Remarks}

\numberwithin{equation}{subsection}

% Macros
%
\def\lim{\mathop{\mathrm{lim}}\nolimits}
\def\colim{\mathop{\mathrm{colim}}\nolimits}
\def\Spec{\mathop{\mathrm{Spec}}}
\def\Hom{\mathop{\mathrm{Hom}}\nolimits}
\def\Ext{\mathop{\mathrm{Ext}}\nolimits}
\def\SheafHom{\mathop{\mathcal{H}\!\mathit{om}}\nolimits}
\def\SheafExt{\mathop{\mathcal{E}\!\mathit{xt}}\nolimits}
\def\Sch{\mathit{Sch}}
\def\Mor{\mathop{\mathrm{Mor}}\nolimits}
\def\Ob{\mathop{\mathrm{Ob}}\nolimits}
\def\Sh{\mathop{\mathit{Sh}}\nolimits}
\def\NL{\mathop{N\!L}\nolimits}
\def\CH{\mathop{\mathrm{CH}}\nolimits}
\def\proetale{{pro\text{-}\acute{e}tale}}
\def\etale{{\acute{e}tale}}
\def\QCoh{\mathit{QCoh}}
\def\Ker{\mathop{\mathrm{Ker}}}
\def\Im{\mathop{\mathrm{Im}}}
\def\Coker{\mathop{\mathrm{Coker}}}
\def\Coim{\mathop{\mathrm{Coim}}}

% Boxtimes
%
\DeclareMathSymbol{\boxtimes}{\mathbin}{AMSa}{"02}

%
% Macros for moduli stacks/spaces
%
\def\QCohstack{\mathcal{QC}\!\mathit{oh}}
\def\Cohstack{\mathcal{C}\!\mathit{oh}}
\def\Spacesstack{\mathcal{S}\!\mathit{paces}}
\def\Quotfunctor{\mathrm{Quot}}
\def\Hilbfunctor{\mathrm{Hilb}}
\def\Curvesstack{\mathcal{C}\!\mathit{urves}}
\def\Polarizedstack{\mathcal{P}\!\mathit{olarized}}
\def\Complexesstack{\mathcal{C}\!\mathit{omplexes}}
% \Pic is the operator that assigns to X its picard group, usage \Pic(X)
% \Picardstack_{X/B} denotes the Picard stack of X over B
% \Picardfunctor_{X/B} denotes the Picard functor of X over B
\def\Pic{\mathop{\mathrm{Pic}}\nolimits}
\def\Picardstack{\mathcal{P}\!\mathit{ic}}
\def\Picardfunctor{\mathrm{Pic}}
\def\Deformationcategory{\mathcal{D}\!\mathit{ef}}

% --- appended 2026-10-12 (Chief): the project own definitions of the reference/slogan/history
% display environments, copied mechanically from the shared Stacks preamble lines 8-12
% (\usepackage{verbatim} plus the three \newenvironment lines, which the chapter preamble keeps
% verbatim from the project preamble). Nothing added, removed or redefined: these environments
% carry no mathematical content and the project itself typesets them as comments. Without these
% three lines a counted statement whose own text uses them fails with Environment slogan undefined.
\usepackage{verbatim}
\newenvironment{reference}{\comment}{\endcomment}
\newenvironment{slogan}{\comment}{\endcomment}
\newenvironment{history}{\comment}{\endcomment}
\usepackage{cleveref}
\crefname{originallemma}{original lemma}{original lemmas}
\Crefname{originallemma}{Original lemma}{Original lemmas}
\makeatletter
\newlabel{lemma-geometrically-normal-formal-fibres-number-of-branches}{{0C2E}{source}{}{sourcepointer-1}{}}
\newlabel{lemma-geometrically-normal-formal-fibres-number-of-branches@cref}{{[originallemma][0][]0C2E}{[1][0][]source}{}{}{}}
\makeatother
\crefname{originallemma}{original lemma}{original lemmas}
\Crefname{originallemma}{Original lemma}{Original lemmas}
\makeatletter
\newlabel{lemma-equal-nr-branches-completion}{{0C29}{source}{}{sourcepointer-2}{}}
\newlabel{lemma-equal-nr-branches-completion@cref}{{[originallemma][0][]0C29}{[1][0][]source}{}{}{}}
\makeatother
\crefname{originalproposition}{original proposition}{original propositions}
\Crefname{originalproposition}{Original proposition}{Original propositions}
\makeatletter
\newlabel{proposition-ratliff}{{0AW6}{source}{}{sourcepointer-3}{}}
\newlabel{proposition-ratliff@cref}{{[originalproposition][0][]0AW6}{[1][0][]source}{}{}{}}
\makeatother
\title{A local Noetherian ring with geometrically normal formal fibres has universally catenary henselization and is itself universally catenary when unibranch}
\author{The Stacks Project Authors}
\date{stacks-project:more-algebra.tex:local-snapshot-20261012; GNU FDL 1.2 or later}
\begin{document}
\maketitle
\setcounter{subsection}{1}
\noindent\textit{Editorial note.} All delivered mathematical text is the author's own, copied from the pinned Stacks Project chapter file more-algebra.tex; no statement or proof step was rewritten, reordered or removed. The counted statement is the labeled lemma environment at source lines 32497-32509 (canonical tag 0C2F) and the counted proof is the proof environment at source lines 32511-32523, which immediately follows it and is its own proof.\par\medskip
\noindent\textit{Editorial note.} The counted statement contains the author's reference environment recording the source of the result; the retained preamble defines that environment as comment-type, so its body and the citation inside it are discarded when the document is typeset, which is why the card needs no bibliography.\par\medskip
\noindent\textit{Editorial note.} The three references are all to earlier results of this chapter, namely the two branch-count statements and the Ratliff criterion that is itself the neighbouring published item; the kinds recorded are the type words appearing in the labels themselves.\par\medskip
\noindent\textit{Editorial note.} The preamble span runs to line 90, the current end of the retained preamble, which is exactly what makes the reference environment a comment; the excerpt serves only as preamble, never as a span coordinate.\par\medskip
\noindent\textit{Editorial note.} The source URL is the source-level Stacks Project URL, as required when more than one card is produced from one source; the per-claim tag is carried by claim\_key and by this note.\par\medskip
\noindent\textit{Allowed standard prerequisites (editorial).}\begin{itemize}
\item For a Noetherian local ring with geometrically normal formal fibres the number of branches of the ring and of its completion agree. invoked at proof lines 32511-32513. Feeds the branch count into the criterion comparing branches with the completion, which is what produces the uniqueness of the minimal prime.
\item If the number of branches of a local ring and its completion agree, then the completion of the henselization modulo a minimal prime has a unique minimal prime. invoked at proof lines 32513-32517. Gives the formal catenarity of the henselization by definition of that notion.
\item A formally catenary Noetherian local ring is universally catenary, and a unibranch ring has henselization with a unique minimal prime. invoked at proof lines 32517-32522. Turns formal catenarity into universal catenarity in the first case, and supplies the uniqueness of the minimal prime that starts the second case.
\end{itemize}\medskip
\ifdefined\hypertarget\hypertarget{sourcepointer-1}{}\fi
\noindent\textit{Source reference.} Original lemma 0C2E: sources/stacks-more-algebra/pinned-original-source.tex, source lines 32231--32231. The reference is applied at proof lines 32511-32512 to compare the numbers of branches of the ring and its completion; the use made of the comparison is inside the counted span.\par\medskip
\ifdefined\hypertarget\hypertarget{sourcepointer-2}{}\fi
\noindent\textit{Source reference.} Original lemma 0C29: sources/stacks-more-algebra/pinned-original-source.tex, source lines 31964--31964. The reference is applied at proof lines 32513-32515 to obtain the uniqueness of the minimal prime of the completion of the henselization; the conclusions drawn from it are inside the counted span.\par\medskip
\ifdefined\hypertarget\hypertarget{sourcepointer-3}{}\fi
\noindent\textit{Source reference.} Original proposition 0AW6: sources/stacks-more-algebra/pinned-original-source.tex, source lines 32484--32484. The reference is invoked twice, at proof lines 32517-32519 and again at proof lines 32522-32523, to pass from formal to universal catenarity; both uses are the last step of their respective case and the preparation is inside the counted span.\par\medskip
\begin{lemma}
\label{lemma-geometrically-normal-fibres-universally-catenary}
\begin{reference}
\cite[Corollary 2.3]{Heinzer-Rotthaus-Wiegand}
\end{reference}
Let $(A, \mathfrak m)$ be a Noetherian local ring with
geometrically normal formal fibres. Then
\begin{enumerate}
\item $A^h$ is universally catenary, and
\item if $A$ is unibranch (for example normal), then
$A$ is universally catenary.
\end{enumerate}
\end{lemma}

\begin{proof}
By Lemma \ref{lemma-geometrically-normal-formal-fibres-number-of-branches}
the number of branches of $A$ and $A^\wedge$ are the same,
hence Lemma \ref{lemma-equal-nr-branches-completion} applies.
Then for any minimal prime $\mathfrak q \subset A^h$
we see that $A^\wedge/\mathfrak q A^\wedge$
has a unique minimal prime. Thus $A^h$ is formally catenary
(by definition) and hence universally catenary by
Proposition \ref{proposition-ratliff}. If $A$ is unibranch,
then $A^h$ has a unique minimal prime, hence $A^\wedge$ has
a unique minimal prime, hence $A$ is formally catenary and
we conclude in the same way.
\end{proof}

\end{document}
